\section{Related Work}
\label{sec:formatting}



\textbf{World Foundation Models}.
As a significant advancement of spatial intelligence, world foundation models (WFMs)~\cite{brooks2024video,deepmind_genie3_2024,ali2025world,meng2024towards,kang2024far} involve the perception, simulation, and interaction within dynamic scenes. Given these properties, 3D geometric perception (covering depth estimation~\cite{piccinelli2024unidepth,yang2024depth,wang2025moge}, scene reconstruction~\cite{wang2025vggt,zhang2024monst3r, wang2025continuous, zhang2025flare,li2025megasam}, and dynamic point tracking~\cite{karaev2025cotracker3,xiao2024spatialtracker,xiao2025spatialtrackerv2}) and video generation (covering text-to-video~\cite{kong2024hunyuanvideo,wang2024scene,hong2022cogvideo,wan2025}, image-to-video~\cite{blattmann2023stable,xing2024dynamicrafter,xu2024easyanimate}, and pose-to-video~\cite{bahmani2025ac3d,he2024cameractrl,bai2025recammaster,bai2024syncammaster}) have emerged as fundamental technologies of WFMs.
This paper presents a unified resource that integrates spatio-temporal semantic and geometric information, advancing WFMs from separate video generation or 3D perception to interactive simulations within virtual environments.



\noindent\textbf{Video Data with Geometric/Semantic Annotations}.
Existing datasets~\cite{yao2020blendedmvs,cabon2020virtual,zhou2018stereo,bai2024syncammaster, zhao2024genxd, rockwell2025dynamic} for 3D perception primarily provide annotations such as depth maps, camera poses, and dynamic tracks, facilitating spatial tasks like depth estimation, scene reconstruction, and dynamic point tracking.
Meanwhile, text-to-video datasets typically consist of video collections with various scales, accompanied by either brief~\cite{soomro2012ucf101,xu_2016_CVPR,zhu2022celebv,yuan2024chronomagic} or detailed~\cite{xue2022advancing,chen2024panda,wang2025koala,nan2024openvid} descriptions.
Despite the availability of these datasets, they frequently lack a comprehensive resource capable of supporting large-scale advancements in both 3D understanding and video generation.
Notably, concurrent studies demonstrate an increasing trend toward integrating spatial geometry and semantic information. However, these works remain constrained either by the data scale (600+ hours of Sekai~\cite{li2025sekai} compared to our 4,000+ hours) or the comprehensive geometric annotations (the lack of consistent 3D point tracks in SpatialVID~\cite{wang2025spatialvid}). As summarized in Table~\ref{tab:source_video}, SceneScribe-1M features comprehensive geometric and semantic annotations for dynamic scenes, demonstrating superior scale and applicability compared to existing datasets.



